\chapter{The forward diffusion}
\label{ch:diffusion}

\section{Destroying data on purpose}

A diffusion model is built around a process that takes a data point and gradually turns it
into noise. The requirements are modest: the process should be easy to simulate, its
end-state should be easy to sample from, and --- crucially --- the distribution at every
intermediate time should be something we can reason about.

The Ornstein--Uhlenbeck (OU) process satisfies all three. Following the convention of the
notes we are matching \citep{mezard2025notes}, it is the stochastic differential equation
\begin{equation}
\label{eq:ou}
\frac{\mathrm{d}x}{\mathrm{d}t} \;=\; -x + \sqrt{2}\,\gwn(t),
\qquad x(0) = a ,
\end{equation}
where $\gwn$ is Gaussian white noise with $\langle \gwn(t)\rangle=0$ and
$\langle \gwn_i(t)\gwn_j(t')\rangle = \delta_{ij}\delta(t-t')$.

Two forces act. The term $-x$ pulls towards the origin; the noise pushes outwards. At long
times they balance and the process forgets where it started.

\section{Solving it}

Equation \eqref{eq:ou} is linear, so it can be integrated explicitly. Multiply by the
integrating factor $e^{t}$:
\[
\frac{\mathrm{d}}{\mathrm{d}t}\!\left(e^{t}x\right)
= e^{t}\!\left(\frac{\mathrm{d}x}{\mathrm{d}t} + x\right)
= \sqrt{2}\,e^{t}\gwn(t).
\]
Integrating from $0$ to $t$ and using $x(0)=a$,
\begin{equation}
x(t) \;=\; e^{-t} a + \sqrt{2}\int_0^t \mathrm{d}\tau\; e^{-(t-\tau)}\gwn(\tau).
\end{equation}
The remaining integral is a sum of independent Gaussian contributions, hence Gaussian, with
mean zero and variance
\[
2\int_0^t \mathrm{d}\tau\; e^{-2(t-\tau)}
= 2\,e^{-2t}\!\int_0^t\! \mathrm{d}\tau\, e^{2\tau}
= 1 - e^{-2t}.
\]
So the whole solution can be written with a single standard Gaussian $z$:
\begin{empheq}[box=\fbox]{equation}
\label{eq:forward-solution}
x(t) \;=\; e^{-t} a + \sqrt{\Dt}\, z ,
\qquad
\Dt \;=\; 1 - e^{-2t},
\qquad z\sim\normal{0}{I}.
\end{empheq}

This one equation is used everywhere below. Read it as: \emph{the signal is scaled down by
$e^{-t}$ and independent Gaussian noise of variance $\Dt$ is added}.

\begin{intuition}
The two coefficients move in opposite directions and are tied together. At $t=0$ we have
$e^{-t}=1$, $\Dt=0$: the data, untouched. As $t\to\infty$, $e^{-t}\to0$ and $\Dt\to1$: pure
standard Gaussian noise, with the data completely gone. The tie
$e^{-2t}+\Dt=1$ means the total variance is preserved when the data has unit variance, so
nothing blows up or collapses along the way.
\end{intuition}

\begin{codebox}
Noising is applied wherever an experiment needs it, always as
\texttt{alpha * a + sqrt(delta) * rng.standard\_normal(...)} with
\texttt{alpha = exp(-t)}, \texttt{delta = 1 - exp(-2t)}. See for instance
\coderef{experiments/exp\_16\_sampling\_validation.py}{fit\_arms}, and
\coderef{src/denoiser.py}{alpha\_delta} for the coefficient pair.
\end{codebox}

\section{The noised distribution \texorpdfstring{$\Pt$}{Pt}}

Equation~\eqref{eq:forward-solution} describes one trajectory. What matters is the
distribution of $x(t)$ when $a$ is drawn from $\Pz$. Since $x$ given $a$ is Gaussian with mean
$e^{-t}a$ and covariance $\Dt I$, marginalising over $a$ gives
\begin{empheq}[box=\fbox]{equation}
\label{eq:pt-def}
\Pt(x) \;=\; \int \mathrm{d}a\; \Pz(a)\,
\frac{1}{(2\pi\Dt)^{\nsites/2}}
\exp\!\left[-\frac{\|x - e^{-t}a\|^2}{2\Dt}\right].
\end{empheq}

That is: $\Pt$ is $\Pz$ convolved with a Gaussian, after a rescaling. Everything about the
forward process is contained in this formula.

\begin{intuition}
Convolution with a Gaussian is smoothing. Whatever sharp features $\Pz$ has --- spikes,
edges, corners --- get blurred, and blurred more as $t$ grows. This is why the problem gets
easier at large $t$ (everything looks Gaussian) and why the interesting behaviour is at small
and intermediate $t$, where the original structure is still partly visible.
\end{intuition}

\section{Why the coordinatewise structure of the noise matters}
\label{sec:coordinatewise}

One feature of \eqref{eq:forward-solution} is so important to this project that it deserves
isolating now, even though its consequence only becomes visible in Chapter~\ref{ch:bp}.

The noise added to coordinate $i$ is independent of the noise added to coordinate $j$. So the
conditional density of the whole observation factorises over sites:
\begin{equation}
\label{eq:likelihood-factorises}
P(x \mid a, t)
= \prod_{i=1}^{\nsites}
\frac{1}{\sqrt{2\pi\Dt}}
\exp\!\left[-\frac{(x_i - e^{-t}a_i)^2}{2\Dt}\right]
\;=\;\prod_{i=1}^{\nsites} \like(x_i \mid a_i).
\end{equation}
Each factor $\like(x_i\mid a_i)$ involves \emph{one} clean coordinate and \emph{one}
observation.

\begin{pitfall}
If the forward noise were correlated across coordinates --- say with covariance $\Sigma_z$
instead of $I$ --- then \eqref{eq:likelihood-factorises} would not factorise, the likelihood
would couple different $a_i$ directly, and the graphical structure exploited in
Chapter~\ref{ch:bp} would be destroyed. Everything in this project depends on the noise being
coordinatewise. It is not a cosmetic modelling choice.
\end{pitfall}

\section{Reversing the process}
\label{sec:reverse}

Generation runs the process backwards: start from noise at a large time and integrate towards
$t=0$. The reverse-time dynamics of \eqref{eq:ou} are
\begin{equation}
\label{eq:reverse}
\frac{\mathrm{d}x}{\mathrm{d}\tau} \;=\; x + 2\,\score(x,\tau) + \gwn(\tau),
\qquad \score(x,t) := \nabla_x \log \Pt(x),
\end{equation}
where $\tau$ runs backwards in $t$. The field $\score$ is the \emph{score}, and it is the only
part of \eqref{eq:reverse} that depends on the data distribution. Everything else is known.

So the entire modelling problem is: \emph{estimate the score}. Chapter~\ref{ch:score} shows
that this is the same as estimating the posterior mean.

\begin{codebox}
\coderef{src/reverse.py}{reverse\_sde} integrates \eqref{eq:reverse} by Euler--Maruyama;
\coderef{src/reverse.py}{probability\_flow\_ode} integrates the deterministic counterpart
with Heun's method. Both accept any score function with signature \texttt{(X, t) -> S}, which
is how the different estimators are compared on equal footing.
\coderef{src/reverse.py}{time\_grid} builds the geometric time discretisation.
\end{codebox}

\begin{pitfall}
Integration cannot be carried to $t=0$. As $\Dt\to0$ the score \eqref{eq:score-final} carries
a factor $1/\Dt$ and diverges. Every implementation stops at some $t_{\min}>0$, and what it
produces is therefore a sample of $P_{t_{\min}}$, \emph{not} of $\Pz$. Chapter~\ref{ch:numerics}
returns to this; it caused two successive errors in this project's own measurements before it
was handled correctly.
\end{pitfall}
